\begin{code}[hide]
module paper.sections.design where

open import Data.Nat using (ℕ ; _+_)
open import Data.Product using (_×_ ; _,_)
open import Data.Sum using (_⊎_)
open import Data.List using (List)
open import Data.Unit using (⊤)
open import Level using (Level ; _⊔_ ; 0ℓ)
open import Relation.Binary using (Rel)
open import Relation.Binary.Definitions using (Decidable)
open import Function.Definitions using (Congruent ; Injective)
open import Algebra.Bundles using (Group)
open import Algebra.Morphism.Structures using (module GroupMorphisms)

open import Notations using (₁₊ ; ₂₊)
open import Word.Base using (wconcatmap)
import Word.Base as WB
open import Presentation.GroupLike using (Grouplike ; module Group-Lemmas)
open import Presentation.Construct.Base using ([_]ₗ ; [_]ᵣ)

postulate
  proof-elided : ∀ {ℓ} {A : Set ℓ} → A
  X Y A B C D  : Set
  a ℓ c d      : Level
  n k          : ℕ
\end{code}
\section{Library Design}\label{sec:design}

This section presents the library bottom-up, with the rationale first.
Five decisions shape everything that follows; each is a bet about where
proof effort should go.
\begin{enumerate}
\item \emph{Syntax is unquotiented, and everything out of it is a function
  on words.}  Words are a plain inductive type and the presented monoid is
  a setoid on them, so normal forms, coset tables, and sections are
  structurally recursive programs that Agda's evaluator can run
  (\S\ref{sec:words}, \S\ref{sec:presentations}).
\item \emph{Specifications are standard-library bundles.}  ``Normal form''
  is not a new notion: it is an injection, a right inverse, or a bijection
  out of the setoid of words (\S\ref{sec:nfproperty}).  No new theory means
  no new lemmas to maintain, and completeness is a five-line lemma.
\item \emph{Certificates are first-order.}  A verified coset table is finite
  data plus five equations quantified over finite types
  (\S\ref{sec:engine}); obligations are therefore finite case splits ---
  which must still be made --- whose cases are closed goals, discharged by
  evaluation where they reduce to a computation and by short equational
  derivations otherwise.
\item \emph{Indices are chosen for the unifier.}  Wire counts are built from
  successors only, so coverage checking succeeds by injectivity of
  constructors rather than by arithmetic (\S\ref{sec:circuits}).
\item \emph{Constructions come in pairs.}  Every product construction has a
  syntactic theorem (normal forms lift) and a semantic one (the constructed
  relation set presents the constructed group), and the semantic groups are
  standard-library bundles (\S\ref{sec:constructions},
  \S\ref{sec:forstdlib}).
\end{enumerate}
Each subsection names its home module in parentheses; code is quoted
verbatim.

\subsection{Words: the free monoid}\label{sec:words}

A word over a generator set $X$ (\texttt{Word.Base}) is an unquotiented
syntax tree --- deliberately \emph{not} a list:

\begin{code}
data Word (X : Set) : Set where
  [_]ʷ : X → Word X
  ε    : Word X
  _•_  : Word X → Word X → Word X
\end{code}
\begin{code}[hide]
infixr 7 _•_
infixl 9 _ʷ _ᵗ
infix 4 _≈_ _===_ _≈ₙ_ _≈₂_ _~_ _===₁_ _===₂_
infix 5 _⋄_⋄_

WRel : Set → Set₁
WRel X = Rel (Word X) 0ℓ

postulate
  w v u w' v' : Word X
  _===_       : WRel X
\end{code}

Keeping composition as a binary node means a word remembers its
bracketing, so terms mirror pen-and-paper calculations and rules may be
applied to any parenthesized subterm; associativity is then one of the
\emph{congruence} generators rather than a property of the carrier, and a
solver discharges re-bracketing obligations (\S\ref{sec:solvers}).  The
functorial and monadic structure is standard: \aid{wmap}, \aid{wconcat},
and the extension of a generator-level substitution to words, written
postfix,

\begin{code}
_ʷ : (X → Word Y) → (Word X → Word Y)
_ʷ = wconcatmap
\end{code}

so \aid{(f ʷ)} is ``the unique monoid homomorphism extending $f$''.  Two
stateful traversals drive everything in \S\ref{sec:engine}: a \emph{coset
action} on letters extends to words by threading the coset left to right,

\begin{code}
_ᵗ : (C → Y → Word X × C) → (C → Word Y → Word X × C)
\end{code}
\begin{code}[hide]
_ᵗ = WB._ᵗ
\end{code}

together with its right-to-left mirror; and a small family of
\emph{conjugation helpers} lifts an action of one alphabet on another
through words, which is what the semidirect-product construction of
\S\ref{sec:constructions} uses to state its word-valued conjugation rule.
A relation set is simply $\aid{WRel}\;X = \aid{Rel}\;(\aid{Word}\;X)\;0\ell$.

\subsection{Presented monoids as inductive congruences}\label{sec:presentations}

Presentations (\texttt{Presentation.Base}, parameterized by
$\Gamma : \aid{WRel}\;X$) follow one two-level convention used everywhere:
\aid{{\_}==={\_}} \emph{always} means the raw axioms ($\Gamma$ itself), and
\aid{{\_}≈{\_}} \emph{always} means the monoid congruence they generate ---
an inductive family with eight constructors:

\begin{code}
data _≈_ : WRel X where
  refl  : w ≈ w
  sym   : w ≈ v → v ≈ w
  trans : w ≈ v → v ≈ u → w ≈ u
  cong  : w ≈ w' → v ≈ v' → w • v ≈ w' • v'
  assoc      : (w • v) • u ≈ w • (v • u)
  left-unit  : ε • w ≈ w
  right-unit : w • ε ≈ w
  axiom      : w === v → w ≈ v
\end{code}

$\approx$ \emph{is} the word problem of $\langle X \mid \Gamma \rangle$,
and a derivation is a first-class syntactic object --- one can compute
with it, induct on it, and translate it (the congruence-lifting lemmas of
\S\ref{sec:constructions} are inductions on this type).  The presented
monoid is thus a setoid on words in the tradition
of~\cite{barthe2003setoids}, not a quotient type as in mathlib's
\aid{PresentedGroup}~\cite{mathlib2020,mathlib-doc-presented} or a higher
inductive type in cubical Agda~\cite{vezzosi2019cubical,cubicallibrary};
\S\ref{sec:setoids} weighs that choice once the case studies are on the
table.
\texttt{Presentation.Properties} packages $\approx$ as an equivalence,
hence as a setoid on words --- the \emph{word setoid}, the domain of every
normal form in \S\ref{sec:nfproperty} --- and as the monoid
$\aid{•-ε-monoid}$, the presented monoid; \texttt{Presentation.GroupLike}
upgrades it to a group when a \aid{Grouplike} witness supplies a left
inverse for every generator: inverses of words are then derived,
cancellation and uniqueness of inverses are lemmas, and
$\aid{•-ε-group}$ is the presented group.

Presentations of concrete algebras are records
(\texttt{Presentation.Definitions}).  The central one:

\begin{code}
record _IsPresentationOf_ (_===_ : WRel X) (G : Group a ℓ) : Set (a ⊔ ℓ) where
  field gl : Grouplike _===_
  module GL = Group-Lemmas _===_ gl
  open GroupMorphisms (Group.rawGroup GL.•-ε-group) (Group.rawGroup G)
  field ⟦_⟧ : Word X → Group.Carrier G
        iso : IsGroupIsomorphism ⟦_⟧
\end{code}
\begin{code}[hide]
postulate
  |NF|       : Set
  nf         : Word X → |NF|
  _≈ₙ_       : |NF| → |NF| → Set
  NormalForm : Set
  inv-nf     : |NF| → Word X
  Sem        : Set
  ⟦_⟧        : Word X → Sem
  _≈₂_       : Sem → Sem → Set
\end{code}

i.e.\ a grouplike witness plus a group isomorphism from the presented
group onto $G$ whose underlying function is the interpretation.  Weaker
records (\aid{{\_}IsSubPresentationOf{\_}} with a monomorphism;
monoid-level variants) come with upgrade lemmas: a surjective
sub-presentation is a presentation, and a grouplike monoid presentation of
$G$'s underlying monoid is a group presentation
(\aid{monoid-presentation⇒group-presentation} in the index) --- the
latter via the \texttt{ForStdlib} fact that monoid homomorphisms between
groups preserve inverses (\S\ref{sec:forstdlib}).  One design point
deserves emphasis: since words have no primitive inverses, \emph{every}
presentation in this library is at bottom a \emph{monoid} presentation,
and group presentations are the \aid{Grouplike} special case.  The
distinction is mathematically real, and the smallest cyclic relation
already exhibits it: read as a group presentation, the trivial relation
$T^{\,0} = \varepsilon$ presents the cyclic group of order $0$ --- by
definition $\mathbb{Z}$ --- but the \emph{monoid} it presents is the
free monoid $(\mathbb{N}, +, 0)$, and grouplikeness is exactly what
fails at order $0$ (\S\ref{sec:cyclic}).  These upgrades structure every proof in
the catalogue: one proves soundness, injectivity (by normalization), and
surjectivity separately, and the records assemble the isomorphism.

\subsection{The zoo of normal-form witnesses}\label{sec:nfproperty}

Normal forms (\texttt{Normalization.NormalForm.Setoid}) are defined
relative to $\Gamma$ and a codomain setoid $\mathit{NF}$, as what it means
for a function out of the word setoid to be a normal form.  The
definitions are, on the nose, the standard library's function bundles
from the word setoid --- no new theory, just names for the roles they
play:
\[
\begin{array}{lcl}
\aid{NormalFormInjective} & = & \aid{Injection}\ \text{(word setoid)}\ \mathit{NF} \\
\aid{NormalForm}          & = & \aid{RightInverse}\ \text{(word setoid)}\ \mathit{NF} \\
\aid{BijectiveNormalForm} & = & \aid{Bijection}\ \text{(word setoid)}\ \mathit{NF}
\end{array}
\]
An \aid{Injection} bundles the map $\aid{nf}$, its congruence, and
injectivity: $w \approx v \iff \aid{nf}\;w \approx_{\mathit{NF}}
\aid{nf}\;v$.  Two consequences ship with it: the workhorse
\begin{code}
by-equal-nf : nf w ≈ₙ nf v → w ≈ v
\end{code}
\begin{code}[hide]
by-equal-nf = proof-elided
\end{code}
which proves word equations by \emph{computing} both normal forms
(\S\ref{sec:solvers}), and decidability transport,
\begin{code}
≈-dec : Decidable _≈ₙ_ → Decidable _≈_
\end{code}
\begin{code}[hide]
≈-dec = proof-elided
\end{code}
so a normal form with decidable carrier \emph{decides the word problem}.
A \aid{RightInverse} additionally carries a section
$\aid{inv-nf} : |\mathit{NF}| \to \aid{Word}\;X$ with
$\aid{inv-nf} \circ \aid{nf} \approx \mathrm{id}$; from it we
derive $\aid{normalize} = \aid{inv-nf} \circ \aid{nf}$, its idempotence,
and injectivity of \aid{nf} --- so every \aid{NormalForm} is a
\aid{NormalFormInjective}.  The bijective witness is a \aid{Bijection}
rather than the standard library's \aid{Inverse}, although the two are
interderivable there, and the choice is one of style with consequences
downstream.  An \aid{Inverse} carries an explicit inverse map together
with \emph{both} inverse laws; what our downstream code consumes is a
section and the single law $\aid{inv-nf} \circ \aid{nf} \approx
\mathrm{id}$, i.e.\ a \aid{RightInverse}, and the second law ---
\emph{exactness}, $\aid{nf} \circ \aid{inv-nf} \equiv \mathrm{id}$ ---
is a hypothesis we deliberately keep separate, because carriers are
allowed to contain junk that no word reaches (a coset type before
pruning, a product carrier before a quotient) and because exactness is
what \aid{by-completeness} below trades against uniqueness.
\aid{BijectiveNormalForm} exists for the cases where surjectivity onto
the carrier is established as a \emph{property} and the section is then
derived from it (as the first projection of the surjectivity witness),
which is how the extension construction of \S\ref{sec:constructions}
consumes its factor normal forms; an \aid{Inverse} would have forced an
explicit inverse map and exactness on every such client.

The converse direction is where constructivity bites, a point we were
careful about.  Classically, any injective $\aid{nf}$ whose image can
be identified has a section by choosing preimages; constructively,
\emph{the section is data}.  An \aid{Injection} gives no algorithm
producing, for a normal form $u$, a word realizing it --- that algorithm
is exactly the extra field of \aid{RightInverse}, and it is what the
completeness lemma below consumes ($\aid{inv-nf}$ realizes each normal
form as a canonical representative).  So the two witnesses are genuinely
different strengths in Agda, injectivity strictly weaker as \emph{data},
even though they classically coincide; we therefore keep both,
use the weak one where only word-problem reasoning is needed
(\aid{by-equal-nf}, decidability), and demand the strong one for
completeness and presentations.  A fourth, deliberately crude witness,
\aid{WeakNormalForm}, drops even the congruence of the map ---
an injective \emph{invariant} suffices to separate words, which is
occasionally all one needs when transporting along unions of relation
sets (\S\ref{sec:constructions}).

Completeness itself is one generic lemma.  Fix a semantics
$\aid{⟦\_⟧} : \aid{Word}\;X \to \mathit{Sem}$ into any setoid.  A
\aid{NormalForm} is \emph{unique for} $\aid{⟦\_⟧}$ when representatives
with equal denotations are equal normal forms:

\begin{code}
record UniqueNormalForm (normalForm : NormalForm) : Set (c ⊔ d) where
  field unique : ∀ {u v : |NF|} → ⟦ inv-nf u ⟧ ≈₂ ⟦ inv-nf v ⟧ → u ≈ₙ v
by-normalization : {normalForm : NormalForm} → UniqueNormalForm normalForm →
  Congruent _≈_ _≈₂_ ⟦_⟧ → Injective _≈_ _≈₂_ ⟦_⟧
\end{code}
\begin{code}[hide]
by-normalization = proof-elided
\end{code}

\emph{Soundness plus a unique normal form yields completeness}: given
$\aid{⟦}w\aid{⟧} \approx_2 \aid{⟦}v\aid{⟧}$, soundness transports along
$w \approx \aid{normalize}\;w$, uniqueness identifies the normal forms,
and injectivity of \aid{nf} closes the loop.  The proof is five lines; its
value is that every case study only ever proves \aid{unique} --- a
statement about \emph{normal forms}, i.e.\ about concrete first-order data
--- never about arbitrary words.  A companion lemma
(\aid{by-normalization-wsurj}) similarly reduces surjectivity of the
semantics to surjectivity on normal forms, and a converse
(\aid{by-completeness}) recovers \aid{UniqueNormalForm} from
completeness plus an exact section ($\aid{nf} \circ \aid{inv\text{-}nf}
\equiv \mathrm{id}$) --- the route by which the construction theorems of
\S\ref{sec:constructions} lift uniqueness witnesses through products,
semidirect products, and amalgamations
(Table~\ref{tab:nfs}).  Sub-presentations then come
for free (\texttt{Normalization.StarPresentation}): given soundness on
axioms, a \aid{NormalForm}, and a \aid{UniqueNormalForm}, the module
assembles the monoid (or, with \aid{Grouplike}, group) monomorphism ---
the injectivity field being literally \aid{by-normalization}.

\subsection{Inductively defined circuits}\label{sec:circuits}

The circuit layer (\texttt{Circuit.Base}) is parameterized by a gate
family $\aid{Gate} : \mathbb{N} \to \mathsf{Set}$ and defines the
generators \aid{Gen} and circuits $\aid{Circuit}\;n =
\aid{Word}\;(\aid{Gen}\;n)$ exactly as in \S\ref{sec:permutations}.  Two
design choices matter more than they look.  First, the indices are
engineered for the unifier.  Addition on naturals recurses on its first
argument, so $k + n$ computes as soon as $k$ is a successor pattern while
$n + k$ is stuck; given that $+$ is left-biased, we make our indices
left-biased too: \aid{gate₂} lands in $\aid{Gen}\;(2 + n)$, not in
$\aid{Gen}\;k$ with a side condition $k \ge 2$, and the $k$-fold shift
has type $\aid{Gen}\;n \to \aid{Gen}\;(k + n)$.  Unification is then all
that is needed to verify the well-formedness constraints that arise ---
every index equation the coverage checker meets is solved by injectivity
of \aid{₁₊}, never by arithmetic, avoiding stuck Diophantine equations of
the form $\mathit{arity} + k \doteq \mathit{target}$ --- and downstream it
is why the coset tables of the
case studies can be defined by direct pattern matching and verified by
\aid{refl}.  Second, the structural rules common to all circuit theories
are factored out once, by \aid{Lift-Relation}, whose permutation
instance appeared in \S\ref{sec:permutations}; its lemma
\aid{lemma-cong↑} lifts the \emph{entire congruence} one wire up ---
proved once, by induction over the eight constructors of $\approx$ ---
so equational developments move freely along the chain
$\mathsf{Cir}\,k \le \mathsf{Cir}\,(k{+}1)$ of \S\ref{sec:circuit-chain}.  A new circuit theory is thus a gate family
\aid{Gate} together with its gate-specific relation set; the structural
rules and their lemmas come for free.

\subsection{The Reidemeister--Schreier engine}\label{sec:engine}

The module \path{Normalization.Reidemeister-Schreier} provides
injectivity theorems for the homomorphisms $(f\,\aid{ʷ})$ induced by
generator maps $f : X \to \aid{Word}\;Y$, in two strengths.
The \emph{simplified} form needs no cosets.  Let $\Gamma$ be a relation
set over $X$ and $\Delta$ one over $Y$, with congruences $\approx_\Gamma$
and $\approx_\Delta$.  If $f$ has a generator-level retraction $g : Y \to
\aid{Word}\;X$ --- $(g\,\aid{ʷ})$ sends each axiom of $\Delta$ to a
$\approx_\Gamma$-equation, and $[x]\aid{ʷ} \approx_\Gamma
(g\,\aid{ʷ})(f\,x)$ for every $x$ --- then $(f\,\aid{ʷ})$ is injective
and $(g\,\aid{ʷ})$ surjective.  This is the tool for
\emph{isomorphisms} of presentations, and powers the final step of each
amalgamation case study via the \aid{StarIsomorphism} builder of
\texttt{Presentation.\allowbreak Morphism}, which turns the four
generator-level conditions into an \aid{IsMonoidIsomorphism} of the
presented monoids.

The \emph{full} form is coset enumeration.  Its packaging
(\texttt{Normalization.CosetNF}) is the heart of the library
(\aid{SingleLevel} shown; a setoid-coset variant exists):

\begin{code}
module SingleLevel (Γ : WRel X) (Δ : WRel Y) (C : Set) (I : C)
  (f : X → Word Y) (h : C → Y → Word X × C) ([_] : C → Word Y)
\end{code}
\begin{code}[hide]
  where

postulate
  I      : C
  f      : X → Word Y
  h      : C → Y → Word X × C
  [_]    : C → Word Y
  _~_    : Word X × C → Word X × C → Set
  _===₁_ : WRel X
  _===₂_ : WRel Y
\end{code}

Think of $\Gamma$ as presenting a subgroup $H$ of the group $G$ presented
by $\Delta$, over the letters $X$ and $Y$ respectively; $C$ is the set of
right cosets, $I$ the identity coset, $f$ the embedding of generators, and
$[\,\_\,]$ the Schreier section.  The table $h$ pushes one letter of $G$
past a coset, emitting a word of $H$ and the updated coset --- so
$\aid{nf} = (h\,\aid{ᵗ})\,I : \aid{Word}\;Y \to \aid{Word}\;X \times C$
runs the whole word from the identity coset.  A submodule
\aid{Transfer} takes the five hypotheses that make the data a genuine
coset presentation:

\begin{code}
module Transfer
  (h=⁻¹f-gen : ∀ (x : X) → ([ x ]ʷ , I) ~ ((h ᵗ) I (f x)))                         -- 1
  (h-wd-ax : ∀ (c : C){u t : Word Y} → u ===₂ t → ((h ᵗ) c u) ~ ((h ᵗ) c t))       -- 2
  (f-wd-ax : ∀ {w v} → w ===₁ v → (f ʷ) w ≈₂ (f ʷ) v)                              -- 3
  ([I]≈ε : [ I ] ≈₂ ε)                                                             -- 4
  (h=ract : ∀ c b → let (b' , c') = h c b in [ c ] • [ b ]ʷ ≈₂ (f ʷ) b' • [ c' ])  -- 5
\end{code}
\begin{code}[hide]
  where
\end{code}

\noindent (1) the table inverts the embedding at the identity coset; (2) the
table respects $\Delta$'s axioms, coset by coset; (3) the embedding
respects $\Gamma$'s axioms; (4) the identity coset is represented by
$\varepsilon$; (5) sliding one letter past a coset's representative
matches the table (the law \aid{ract-sound} of \S\ref{sec:permutations}).
From these, \aid{Transfer} derives well-definedness of
\aid{nf} on $\approx_2$, the factorization
$[c] \aidop{•} b \approx_2 (f\,\aid{ʷ})\,b' \aidop{•} [c']$ for whole words,
injectivity --- and the payoff, the \emph{normal-form transport}
\[
  \aid{nf-transp} : \mathrm{NormalForm}\;\Gamma\;\mathit{NF}
    \to \mathrm{NormalForm}\;\Delta\;(\mathit{NF} \times C),
\]
a normal form for the subgroup yields one for the group, with one more
coset digit.  A record \aid{PackedCosetTable} bundles the data and
hypotheses as a single value --- so a whole verified level is an ordinary
first-class object, two of which make an amalgamation
(\S\ref{sec:constructions}) --- and \aid{CosetTower} iterates levels up an
$\mathbb{N}$-indexed family of presentations, one such record
(\aid{Extension}) per level, with carrier
$\aid{tower-carrier}\;\mathit{NF}_0\;(k{+}1) =
 \aid{tower-carrier}\;\mathit{NF}_0\;k \times C_k$: the mixed-radix normal
form of \S\ref{sec:preliminaries}, one verified digit per level.

Two facts make this engine pleasant in practice.  The hypotheses are
\emph{quantifier-small}: each ranges over cosets and generators (finite
data types) and axioms (a finite inductive family), so in every case
study they are finite case splits.  And once a normal form for $\Gamma$
is in scope, the cases of an obligation like (2) close by
\emph{computation}: each becomes \aid{by-equal-nf Eq.refl} --- both
sides normalize to the same value, and Agda's evaluator checks it.  In
the qutrit Clifford+$T$ tower, a nine-coset level imposes seventy-two
such obligations; every one is discharged by \aid{refl}.  The other
hypotheses --- the soundness law (5) above all --- are equational
derivations done by hand, one per case, as \aid{ract-sound} was in
\S\ref{sec:permutations}.  This is the sense in which the framework
absorbs part of the ``gigantic computations'' of \S\ref{sec:intro}: the
case splits remain, but the cases that are mere computation cost
nothing.

\subsection{Products of presentations, and their presentation theorems}\label{sec:constructions}

Relation sets over a disjoint union of alphabets
(\texttt{Presentation.Construct.Base}) are built from a single primitive,
a three-way join whose \emph{mixed} component is the design parameter:

\begin{code}
data _⋄_⋄_ (Γ₁ : WRel A) (Γ₂ : WRel B) (Γ₃ : WRel (A ⊎ B)) : WRel (A ⊎ B) where
  left  : Γ₁ u v → (Γ₁ ⋄ Γ₂ ⋄ Γ₃) [ u ]ₗ [ v ]ₗ
  right : Γ₂ u v → (Γ₁ ⋄ Γ₂ ⋄ Γ₃) [ u ]ᵣ [ v ]ᵣ
  mid   : Γ₃ u v → (Γ₁ ⋄ Γ₂ ⋄ Γ₃) u v
\end{code}

Choosing the mixed relation recovers the classical constructions ---
free product ($\aid{EmptyRel}$), direct product ($\aid{CommRel}$:
generators of the two sides commute), semidirect product
($\aid{ConjRel}$/$\aid{ConjRelʷ}$: moving $h$ past $n$ conjugates $n$,
by a generator or by a word), amalgamated free product
($\aid{AmalgRel}\;f_1\,f_2$: the two images of a shared alphabet are
identified), and definitional extension ($\aid{SugarRel}$: a new
generator desugars to a word) --- plus $n$-fold iterations
$\Gamma \mathrel{\aid{⊕\char94}} n$.  Congruence-lifting lemmas (\aid{lefts},
\aid{rights}) embed the factors' congruences into the join, and
normal-form transports move witnesses along unions and monomorphisms.

For each construction, \texttt{Presentation.Construct.Properties} proves
two kinds of theorems.  \emph{Normal-form lifting}: from NF witnesses for
$\Gamma$ and $\Delta$, an NF witness for the product relation set, with
product carrier ($\mathit{NF}_1 \times \mathit{NF}_2$ for direct and
semidirect products; for amalgamations the classical alternating carrier
\begin{code}[hide]
amalNFC : Set → Set → Set
\end{code}
\begin{code}
amalNFC C D = (D ⊎ ⊤) × List (C × D) × (C ⊎ ⊤)
\end{code}
--- a word of alternating nontrivial coset representatives between two
tails~\cite{magnus2004combinatorial}).  \emph{Presentation}: if the factors
present concrete groups, the constructed relation set presents the
constructed group.  With their premise telescopes elided, the direct,
semidirect, $n$-fold, and amalgamated statements read:
\[
\begin{array}{l}
\aid{dpres} : (\Gamma \mathrel{\aidop{⋄}} \Delta \mathrel{\aidop{⋄}} \aid{CommRel})
   \ \aid{IsPresentationOf}\ (G_1 \times G_2) \\[3pt]
\aid{dpres} : (\Gamma \mathrel{\aidop{⋄}} \Delta \mathrel{\aidop{⋄}} \aid{ConjRelʷ}\;\mathit{conj})
   \ \aid{IsPresentationOf}\ (G_1 \rtimes G_2) \\[3pt]
\aid{presentation} : \forall n \to (\Gamma \mathrel{\aid{⊕\char94}} n)
   \ \aid{IsPresentationOf}\ (\otimes\aid{-group}\;n) \\[3pt]
\aid{dpres} : (P_1 \mathrel{\aid{*}} P_2 \mathrel{\aid{⋆}} f_1 \mathrel{\aid{⋆}} f_2)
   \ \aid{IsPresentationOf}\ (G_1 \ast_{G_0} G_2)
\end{array}
\]
where the semantic groups on the right are built by the
\texttt{ForStdlib} constructions below.  For the amalgamated product the
elided premises are two \aid{PackedCosetTable}s over a common base
presentation (the record \aid{AmalDataNF}) and presentations of the three
groups; the amalgamating homomorphisms $\varphi, \psi$ are then
\emph{derived} from the syntactic embeddings through the presentation
isomorphisms.  A further module proves the classical recipe for
presenting a group \emph{extension}: if $1 \to N \to G \to Q \to 1$ is a
short exact sequence, and presentations of $N$ and $Q$ are given together
with the data that realizes it in $G$ --- how $G$ conjugates $N$, and the
correction in $N$ by which each lifted relator of $Q$ fails to hold ---
then $N$'s relations, word-valued conjugation rules, and the corrected
lifts of $Q$'s relators together present $G$ (cf.\ Proposition~2.55
of~\cite{holt2005handbook}); semidirect products are the special case
with trivial corrections.  These construction theorems
are what make the catalogue of \S\ref{sec:examples} cheap: the Pauli
presentation, for instance, is \emph{assembled} --- cyclic factors, one
binary product, one $n$-fold product --- with no fresh coset enumeration
at all.

\subsection{The semantic side: contributions to the standard library}\label{sec:forstdlib}

The right-hand sides of the presentation theorems need the constructed
groups to \emph{exist} as standard-library-style bundles, and the
standard library did not have them.  The library ships them in a
\texttt{ForStdlib} hierarchy that mirrors stdlib paths for upstreaming.
The module \path{Algebra.Construct.SemiDirectProduct} builds $N \rtimes H$
from raw monoids up to groups, driven by a bundled \aid{Action} record
whose six fields (the action, its congruence, and its preservation of
units and products on both sides) are exactly the laws that make the
twisted multiplication associative and unital.
\path{Algebra.Construct.Amalgamation} builds the amalgamated free product
$M \ast_C N$ of \emph{monoids} with no concrete carrier: elements are
words over $|M| \uplus |N|$ and equality is the least congruence
multiplying out same-factor letters, deleting units, and gluing
$\varphi\,c$ to $\psi\,c$ --- the pushout in the category of monoids when
$\varphi, \psi$ are homomorphisms.  \path{Algebra.Construct.Extension}
and \path{.SplitExtension} are short exact sequences of groups without
and with a homomorphic section, semidirect products being the canonical
split model.  \path{Algebra.Morphism.Consequences} proves that a monoid
homomorphism between groups is a group homomorphism (previously derivable
only via a deprecated stdlib module), and
\path{Algebra.Morphism.Structures} adds \aid{IsMonoidEpimorphism} (stdlib
has monomorphisms and isomorphisms, but no epimorphisms).  Finally,
\path{Data.Fin.Permutation.Properties} bundles stdlib's
\aid{Permutation′}~$n$ with composition, identity, and \aid{flip} as the
\aid{Group} used by the tight semantics of \S\ref{sec:permutations}.  None
of these is deep --- the seven modules total 865 lines --- but each was
necessary for some semantic side of a presentation theorem.  The moral
is that the standard library's algebra hierarchy is rich in
\emph{structures} but thin in \emph{constructions}, and
presentation-level results force the constructions into existence.

\subsection{Proof engineering: solvers and evaluation}\label{sec:solvers}

Three small tools carry a large share of the equational text.
\aid{by-assoc} proves $w \approx v$ whenever the two words have the same
letter sequence --- it normalizes both sides to lists and compares them
with \aid{refl} --- so all re-bracketing and unit shuffling in a
calculation is one appeal to the solver; the pattern-guided variant
\aid{by-passoc} re-brackets toward a shape specified by a pattern word
with holes, useful when the \emph{next} rule application needs a
particular association.  \aid{by-equal-nf} (\S\ref{sec:nfproperty})
closes any goal whose two sides have the same computed normal form; since
coset tables and normal-form maps are structurally recursive programs,
this turns the well-definedness half of table verification into
evaluation.  Both are used like tactics, but neither uses
metaprogramming: each is a verified function into a canonical data type
followed by \aid{refl}, the small-scale-reflection recipe with the
reflection left to the evaluator.
Finally, mundane but effective: a shared \texttt{Notations} module of
numeral patterns (\aid{₀} … \aid{₁₅}, successor patterns \aid{₁₊} …
\aid{₄₊} overloaded over $\mathbb{N}$ and $\mathsf{Fin}$, and \aid{auto}
for \aid{refl}) keeps indices readable, and equational proofs are written
in the standard library's setoid-reasoning style, so the Agda text of,
say, \aid{lemma-XS} in the Clifford+$T$ study reads line for line like
the pen-and-paper derivation it replaces.
